\documentclass[10pt,a4paper]{article}
\usepackage[top=1.5cm, bottom=1.5cm, left=1.6cm, right=1.6cm, headheight=14pt, footskip=18pt]{geometry}
\usepackage[utf8]{inputenc}
\usepackage{amsmath,amsfonts,amssymb}
\usepackage{graphicx}
\usepackage{tikz}
\usepackage{circuitikz}
\usetikzlibrary{calc,patterns}
\usepackage[most]{tcolorbox}
\usepackage{enumitem}
\usepackage{fancyhdr}
\usepackage{booktabs}
\usepackage{tabularx}
\usepackage{colortbl}
\usepackage{hyperref}

% --- Palet Warna Resmi UNIB ---
\definecolor{unibblue}{RGB}{0, 32, 96}
\definecolor{unibgold}{RGB}{197, 160, 89}
\definecolor{unibgreen}{RGB}{34, 112, 60}
\definecolor{boxbg}{RGB}{248, 250, 253}
\definecolor{linegray}{RGB}{210, 215, 225}

% --- Pengaturan Header & Footer ---
\pagestyle{fancy}
\fancyhf{}
\lhead{\footnotesize\textbf{\color{unibblue}Universitas Bengkulu} \textbar\ Program Studi S1 Teknik Elektro}
\rhead{\footnotesize\color{gray}Persamaan Diferensial (Genap 2026)}
\lfoot{\footnotesize\color{gray}LKM OBE \textbar\ \textbf{Video}: \url{https://youtu.be/x15SuJBvZQE} \textbar\ \href{https://www.youtube.com/playlist?list=PLS5oOZWXZeTw}{\color{unibblue}\textbf{Playlist}}}
\cfoot{\footnotesize\color{gray}Hal. \thepage\ dari 2}
\rfoot{\footnotesize\color{unibblue}\textbf{Minggu 5: PDB Orde 2 Non-Homogen \& Resonan...}}
\renewcommand{\headrulewidth}{0.6pt}
\renewcommand{\footrulewidth}{0.4pt}

% --- Custom Box untuk Lembar Jawab ---
\newtcolorbox{answerbox}[1]{%
  colback=white,
  colframe=unibblue!70!black,
  coltitle=white,
  fonttitle=\bfseries\scriptsize,
  colbacktitle=unibblue,
  title={#1},
  arc=2pt,
  left=5pt,right=5pt,top=3pt,bottom=3pt,
  boxrule=0.6pt
}

\begin{document}

% ==================== HALAMAN 1 ====================
\noindent\begin{minipage}[c]{0.68\textwidth}%
  {\large\textbf{\color{unibblue}LEMBAR KERJA MAHASISWA (LKM)}}\\[1mm]
  {\normalsize\textbf{Mata Kuliah: Persamaan Diferensial (2 SKS)}}\\[0.5mm]
  {\small\textbf{Topik Minggu 5:} PDB Orde 2 Non-Homogen \& Resonansi RLC Eksitasi AC}%
\end{minipage}%
\hfill%
\begin{minipage}[c]{0.30\textwidth}%
  \begin{tcolorbox}[colback=unibblue!5,colframe=unibblue,boxrule=0.6pt,arc=2pt,left=4pt,right=4pt,top=2pt,bottom=2pt]
    \scriptsize
    \textbf{Nama:} \dotfill\\[1.2mm]
    \textbf{NPM:} \dotfill\\[1.2mm]
    \textbf{Kelas/Klp:} \dotfill\\[1.2mm]
    \textbf{Nilai Final:} \hfill \textbf{/ 100}
  \end{tcolorbox}%
\end{minipage}\par

\vspace{1.5mm}

% Kotak Sub-CPMK & Pedoman
\begin{tcolorbox}[colback=boxbg,colframe=unibgold!90!black,boxrule=0.6pt,arc=2pt,left=5pt,right=5pt,top=3pt,bottom=3pt,title=\textbf{\scriptsize Capaian Pembelajaran (Sub-CPMK 3) \& Pedoman Evaluasi OBE}]
  \scriptsize
  \textbf{Sub-CPMK 3:} Mahasiswa mampu memecahkan PDB linier orde 2 non-homogen menggunakan metode koefisien tak tentu dan variasi parameter untuk menganalisis resonansi dan faktor kualitas rangkaian AC.\\
  \textbf{Video Panduan (YouTube):} \url{https://youtu.be/x15SuJBvZQE} \hfill \textbf{Playlist:} \href{https://www.youtube.com/playlist?list=PLS5oOZWXZeTw}{\color{unibblue}\textbf{bit.ly/pd-unib-playlist}}\\\textbf{Alokasi Waktu:} 50--60 Menit \quad\textbar\quad \textbf{Model:} Kolaboratif / Mandiri Terstruktur \quad\textbar\quad \textbf{Bahasa Komputasi:} Julia 1.12+
\end{tcolorbox}

\vspace{1mm}

% Soal C1
\noindent\textbf{1. (C1 - Mengingat \textbar\ Bobot: 10 Poin)}\par\vspace{0.5mm}
{\footnotesize Jelaskan prinsip superposisi solusi total PDB non-homogen $y(t) = y_h(t) + y_p(t)$, dan sebutkan aturan pemilihan bentuk tebakan solusi partikular $y_p(t)$ jika fungsi pemaksa berbentuk $F(t) = V_m \sin(\omega t)$.}
\begin{answerbox}{Ruang Pengerjaan C1 (Prinsip Solusi Lengkap \& Bentuk Tebakan Partikular)}
  \vspace{2.0cm}
\end{answerbox}

\vspace{1mm}

% Soal C2
\noindent\textbf{2. (C2 - Memahami \textbar\ Bobot: 15 Poin)}\par\vspace{0.5mm}
{\footnotesize Jelaskan fenomena resonansi listrik seri secara fisis. Mengapa pada frekuensi resonansi $\omega_0 = 1/\sqrt{LC}$, impedansi total rangkaian menjadi minimum ($Z = R$) dan arus mencapai nilai maksimum?}
\begin{answerbox}{Ruang Pengerjaan C2 (Intuisi Fisika Resonansi Seri \& Pembatalan Reaktansi)}
  \vspace{2.1cm}
\end{answerbox}

\vspace{1mm}

% Soal C3
\noindent\textbf{3. (C3 - Menerapkan \textbar\ Bobot: 20 Poin)}\par\vspace{0.5mm}
\noindent\begin{minipage}[t]{0.60\textwidth}%
  {\footnotesize Rangkaian RLC seri di samping dihubungkan ke sumber tegangan AC $v_s(t) = 100 \cos(1000 t)\text{ V}$. Dengan $R = 10\,\Omega$, $L = 0{,}05\text{ H}$, dan $C = 20\,\mu\text{F}$, tentukan solusi partikular tunak (*steady-state*) arus $i_{\text{ss}}(t) = I_m \cos(1000 t - \phi)$.}%
\end{minipage}%
\hfill%
\begin{minipage}[t]{0.37\textwidth}%
  \centering
  \resizebox{0.95\linewidth}{!}{%
  \begin{circuitikz}[american, scale=0.65, transform shape]
    \draw (0,0) to[sV, l=$v_s(t)$] (0,1.8)
          to[R, l=$R$] (1.6,1.8)
          to[L, l=$L$] (3.0,1.8)
          to[C, l=$C$] (3.0,0)
          to[short] (0,0);
    \draw[->, >=stealth, thick, unibblue] (1.4,0.8) arc (160:-160:0.35);
    \node at (1.8,0.8) [unibblue, font=\tiny\bfseries] {$i(t)$};
  \end{circuitikz}%
  }%
\end{minipage}\par\vspace{1mm}
\begin{answerbox}{Ruang Pengerjaan C3 (Penurunan Solusi Partikular Arus AC Tunak)}
  \vspace{2.8cm}
\end{answerbox}

\newpage
% ==================== HALAMAN 2 ====================

% Soal C4
\noindent\textbf{4. (C4 - Menganalisis \textbar\ Bobot: 20 Poin)}\par\vspace{0.5mm}
{\footnotesize Analisis perbesaran tegangan kapasitor pada kondisi resonansi: Definisikan Faktor Kualitas $Q = \frac{\omega_0 L}{R}$. Buktikan secara matematis bahwa pada kondisi resonansi, amplitudo tegangan pada kapasitor memenuhi hubungan $V_{C,\text{maks}} = Q \cdot V_s$. Apa bahaya fenomena ini terhadap dielektrik isolasi?}
\begin{answerbox}{Ruang Pengerjaan C4 (Analisis Faktor Kualitas Q \& Perbesaran Tegangan Kritis)}
  \vspace{3.8cm}
\end{answerbox}

\vspace{1.5mm}

% Soal C5
\noindent\textbf{5. (C5 - Mengevaluasi \textbar\ Bobot: 20 Poin)}\par\vspace{0.5mm}
{\footnotesize Suatu sistem mengalami eksitasi tepat pada frekuensi alami tanpa redaman: $y'' + \omega_0^2 y = F_0 \cos(\omega_0 t)$. Evaluasi kegagalan tebakan standar dan buktikan dengan metode variasi parameter bahwa solusinya tumbuh secara linier terhadap waktu ($t \sin(\omega_0 t)$) menuju resonansi tak-hingga (destruktif).}
\begin{answerbox}{Ruang Pengerjaan C5 (Evaluasi Resonansi Murni Tak-Hingga \& Amplifikasi Linier)}
  \vspace{3.8cm}
\end{answerbox}

\vspace{1.5mm}

% Soal C6
\noindent\textbf{6. (C6 - Merancang \& Komputasi Julia \textbar\ Bobot: 15 Poin)}\par\vspace{0.5mm}
{\footnotesize Rancang program simulasi frekuensi sapuan (\textit{frequency sweep}) dalam \textbf{Julia} untuk menghitung respons amplitudo arus terhadap variasi frekuensi eksitasi $\omega \in [0{,}2\omega_0, 2\omega_0]$. Tuliskan algoritma vektorisasi perhitungan magnitudo impedansi $|Z(\omega)|$.}
\begin{answerbox}{Ruang Pengerjaan C6 (Skrip Julia Kurva Respon Resonansi \& Lebar Pita)}
  \vspace{3.8cm}
\end{answerbox}

\vspace{1.5mm}

% Tabel Rekapitulasi Skor OBE & Tanda Tangan
\noindent\begin{minipage}[b]{0.65\textwidth}%
  \centering
  \scriptsize
  \begin{tabular}{|c|c|c|c|c|c|c|}
    \hline
    \rowcolor{unibblue}
    \textbf{\color{white}C1} & \textbf{\color{white}C2} & \textbf{\color{white}C3} & \textbf{\color{white}C4} & \textbf{\color{white}C5} & \textbf{\color{white}C6} & \textbf{\color{white}Total} \\
    \rowcolor{unibblue!10}
    10 Pts & 15 Pts & 20 Pts & 20 Pts & 20 Pts & 15 Pts & 100 Pts \\
    \hline
    & & & & & & \\[2.5mm]
    \hline
  \end{tabular}%
\end{minipage}%
\hfill%
\begin{minipage}[b]{0.32\textwidth}%
  \centering
  \scriptsize
  \textbf{Verifikasi Dosen / Asisten:}\\[5mm]
  (\dotfill)
\end{minipage}\par

\end{document}
